\input qpd-bootstrap

\begin{problem}{20260727}
  \meta{name}{The Magical Whip}
  \meta{setter}{Andy He}
  \meta{topic}{mechanics}
  \meta{difficulty}{4}
  \meta{tags}{mechanics, variable mass, momentum, Newton's laws}

SM-senior 黄瑞轩 is testing a whip.  But this whip is brittle and magical: it has \emph{no structural integrity}, so each element $\mathrm{d}m$ of it is in free fall on its own.  He holds the whip vertically with its lower end just touching the ground, then lets go.

The whip has mass $M$ and length $L$, with uniform linear density $\lambda = M/L$.  At a later instant a length $y$ of it lies on the ground (piled up at rest), while the remaining length $L - y$ is still falling, the whole falling part descending with speed $v$ as a column.

\begin{assumptions}
  \item The whip is uniform with linear density $\lambda = M/L$.
  \item There is no internal tension: each element is in free fall, and on touching the ground it is instantly brought to rest (perfectly inelastic landing).
  \item The pile on the ground is at rest and does not bounce.
  \item Air resistance is negligible.
\end{assumptions}

\begin{questions}
  \item \textbf{The ground's normal force.}
        Solve for the normal force $F_N$ applied by the ground on the whip as a function of $y$.
\end{questions}

\end{problem}

\begin{solution}{20260727}

This is a \emph{variable-mass} system: the moving part of the whip shrinks as it piles up, so Newton's law in the form $\sum F = m a$ does not directly apply.  Instead we use the momentum form
\[
  \dot{P} = F_{\text{ext}},
\]
where $P$ is the total vertical momentum of the whip (the pile contributes none, being at rest) and $F_{\text{ext}}$ is the total external force on the whip.

The falling column has length $L-y$ and, because every element is in free fall together, a common downward speed $v = \dot{y}$.  Its momentum (downward) is
\[
  P = \lambda\,(L-y)\,v .
\]
Differentiating, with $\dot v = g$ (free fall) and $\dot y = v$,
\[
  \dot{P} = \lambda\,(L-y)\,\dot{v} - \lambda\,\dot{y}\,v
          = \lambda\,(L-y)\,g - \lambda\,v^{2}.
\]

The external forces on the whole whip are its weight $\lambda L g$ (downward) and the ground's normal force $F_N$ (upward).  Taking downward as positive,
\[
  F_{\text{ext}} = \lambda L g - F_N .
\]
Thus $\dot P = F_{\text{ext}}$ gives
\[
  \lambda L g - F_N = \lambda\,(L-y)\,g - \lambda\,v^{2}
  \;\Longrightarrow\;
  F_N = \lambda\,g y + \lambda\,v^{2}.
\]

Finally, the falling column has descended a distance $y$, so by free-fall kinematics
\[
  v^{2} = 2gy .
\]
Substituting,
\[
  \boxed{F_N = 3\lambda g y = \frac{3Mgy}{L}} .
\]

\textbf{Interpretation.}  The term $\lambda g y = Mgy/L$ is the weight of the pile already on the ground; the extra $2\lambda g y$ is the momentum flux — the rate at which the falling whip's momentum is destroyed as each element is brought to rest.  This is why the ground feels \emph{three times} the weight of the already-landed portion.

\end{solution}

\input qpd-end
